\documentclass[a4paper,12pt]{article}
\usepackage[utf8]{inputenc}
\usepackage[T1]{fontenc}
\usepackage[french]{babel}
\usepackage{amsmath, amssymb, amsfonts}
\usepackage{array, longtable, booktabs}
\usepackage[table]{xcolor} % option table pour \rowcolors
\usepackage{geometry}
\geometry{left=1.8cm, right=1.8cm, top=2cm, bottom=2cm}
\usepackage{enumitem}
\usepackage{hyperref}
\hypersetup{colorlinks=true, linkcolor=blue, urlcolor=blue}

% Commandes pour les coches
\newcommand{\fait}{\ensuremath{\checkmark}}
\newcommand{\revoit}{\ensuremath{\circlearrowright}}
\newcommand{\casevide}{\ensuremath{\square}}

% Types de colonnes pour un contrôle précis des largeurs
\newcolumntype{C}[1]{>{\centering\arraybackslash}p{#1}}
\newcolumntype{L}[1]{>{\raggedright\arraybackslash}p{#1}}
\newcolumntype{R}[1]{>{\raggedleft\arraybackslash}p{#1}}

\title{Programmation annuelle de Mathématiques complémentaires -- Terminale générale \\
	Année 2026--2027 -- Zone C (3 h/semaine)}
\author{}
\date{Version du \today}

\begin{document}
	
	\maketitle
	
	\section*{Présentation générale}
	Ce document propose une organisation spiralaire de l’enseignement de l’option Mathématiques complémentaires, conforme au programme officiel. L’année comporte 31 semaines de cours effectifs (hors vacances et jours fériés). Chaque notion est introduite une première fois, puis réinvestie et approfondie dans un second temps, à l’occasion de thèmes d’étude différents. Cette approche favorise une acquisition progressive et durable des connaissances.
	
	\medskip
	\textbf{Place de Python :} La programmation est intégrée de façon transverse et évaluée dans \emph{tous} les sujets (formatives, contrôles, DM, PBL). Les compétences algorithmiques (boucles, tests, fonctions, simulation, calcul approché) sont développées en parallèle des mathématiques et réactivées régulièrement selon le même principe spiralé.
	
	\medskip
	\textbf{Utilisation de l’IA dans les DM :} Dans chaque devoir maison, l’élève est invité à utiliser un outil d’intelligence artificielle générative (ChatGPT, Mistral, Copilot, etc.) pour :
	\begin{itemize}
		\item obtenir des explications ou des exemples supplémentaires sur une notion,
		\item suggérer des pistes de résolution,
		\item débuguer ou améliorer un code Python,
		\item confronter différentes méthodes de résolution.
	\end{itemize}
	\emph{L’IA ne doit cependant pas rédiger la solution finale à la place de l’élève.} Le raisonnement personnel et la rédaction restent exigibles. Un **journal de bord** (questions posées, réponses obtenues, critiques et adaptations apportées) est à annexer à chaque DM. Cette pratique développe l’esprit critique et la capacité à évaluer la pertinence des réponses générées.
	
	\section{Tableau récapitulatif des contenus à enseigner}
	Le tableau ci-dessous reprend l’intégralité des attendus du programme. La colonne \textbf{« Traité »} indique si le contenu a été abordé en classe (coche \fait), et la colonne \textbf{« Reprise »} signale les notions qui feront l’objet d’un second passage (\revoit) pour consolider la compréhension.
	
	\rowcolors{2}{gray!8}{white}
	\begin{longtable}{|L{14cm}|C{1.5cm}|C{1.5cm}|}
		\hline
		\textbf{Thème / Contenu} & \textbf{Traité} & \textbf{Reprise} \\
		\hline
		\endhead
		\hline
		\multicolumn{3}{r}{Suite de la page suivante} \\
		\endfoot
		\hline
		\endlastfoot
		
		% --- Analyse – Suites numériques, modèles discrets ---
		\multicolumn{3}{|L{14cm}|}{\textbf{Analyse – Suites numériques, modèles discrets}} \\
		Approche intuitive de la limite d’une suite (finie, infinie) & \casevide & \casevide \\
		Opérations sur les limites, théorème des gendarmes & \casevide & \casevide \\
		Limite d’une suite géométrique de raison positive & \casevide & \casevide \\
		Limite de la somme des termes d’une suite géométrique (raison $<1$) & \casevide & \casevide \\
		Suites arithmético-géométriques (recherche d’une solution particulière constante) & \casevide & \casevide \\
		Représentation graphique d’une suite récurrente $u_{n+1}=f(u_n)$ & \casevide & \casevide \\
		\hline
		
		% --- Analyse – Fonctions : continuité, dérivabilité, limites ---
		\multicolumn{3}{|L{14cm}|}{\textbf{Fonctions : continuité, dérivabilité, limites}} \\
		Notion de limite, asymptotes horizontales/verticales & \casevide & \casevide \\
		Limites des fonctions de référence (carré, cube, racine, inverse, $\exp$, $\ln$) & \casevide & \casevide \\
		Théorème des valeurs intermédiaires (admis) & \casevide & \casevide \\
		Réciproque d’une fonction continue strictement monotone & \casevide & \casevide \\
		Fonction logarithme népérien (définition, équation fonctionnelle, dérivée) & \casevide & \casevide \\
		Dérivée de $x \mapsto f(ax+b)$, $e^{u(x)}$, $\ln u(x)$, $u(x)^2$ & \casevide & \casevide \\
		\hline
		
		% --- Primitives et équations différentielles ---
		\multicolumn{3}{|L{14cm}|}{\textbf{Primitives et équations différentielles}} \\
		Notion de primitive, lien avec $y'=f$ & \casevide & \casevide \\
		Deux primitives diffèrent d’une constante & \casevide & \casevide \\
		Résolution de $y' = ay + b$ (solution particulière constante, solution générale) & \casevide & \casevide \\
		\hline
		
		% --- Fonctions convexes ---
		\multicolumn{3}{|L{14cm}|}{\textbf{Fonctions convexes}} \\
		Dérivée seconde, définition par la position des sécantes/tangentes & \casevide & \casevide \\
		Caractérisation par $f'' \ge 0$ (admise) & \casevide & \casevide \\
		Point d’inflexion & \casevide & \casevide \\
		\hline
		
		% --- Intégration ---
		\multicolumn{3}{|L{14cm}|}{\textbf{Intégration}} \\
		Intégrale d’une fonction continue positive (aire sous la courbe) & \casevide & \casevide \\
		Relation de Chasles, valeur moyenne & \casevide & \casevide \\
		Approximation par la méthode des rectangles & \casevide & \casevide \\
		Intégrale d’une fonction de signe quelconque & \casevide & \casevide \\
		Théorème fondamental : $F(x)=\int_a^x f(t)\,dt$ est une primitive & \casevide & \casevide \\
		Calcul d’intégrales à l’aide de primitives & \casevide & \casevide \\
		\hline
		
		% --- Probabilités et statistique – Lois discrètes ---
		\multicolumn{3}{|L{14cm}|}{\textbf{Probabilités – Lois discrètes}} \\
		Loi uniforme sur $\{1,\dots,n\}$, espérance & \casevide & \casevide \\
		Épreuve et loi de Bernoulli (espérance, écart-type) & \casevide & \casevide \\
		Schéma de Bernoulli, arbre & \casevide & \casevide \\
		Coefficients binomiaux, triangle de Pascal & \casevide & \casevide \\
		Loi binomiale $B(n,p)$ (espérance, écart-type admis) & \casevide & \casevide \\
		Loi géométrique (définition, espérance admise, absence de mémoire) & \casevide & \casevide \\
		\hline
		
		% --- Lois à densité ---
		\multicolumn{3}{|L{14cm}|}{\textbf{Lois à densité}} \\
		Notion de loi à densité, fonction de répartition & \casevide & \casevide \\
		Espérance et variance d’une loi à densité (intégrales) & \casevide & \casevide \\
		Loi uniforme sur $[0,1]$, puis sur $[a,b]$ & \casevide & \casevide \\
		Loi exponentielle (densité, fonction de répartition, espérance, absence de mémoire) & \casevide & \casevide \\
		\hline
		
		% --- Statistique à deux variables ---
		\multicolumn{3}{|L{14cm}|}{\textbf{Statistique à deux variables}} \\
		Nuage de points, point moyen & \casevide & \casevide \\
		Ajustement affine, droite des moindres carrés, coefficient de corrélation & \casevide & \casevide \\
		Ajustement par changement de variable & \casevide & \casevide \\
		Interpolation / extrapolation & \casevide & \casevide \\
		\hline
		
		% --- Thèmes d’étude transversaux ---
		\multicolumn{3}{|L{14cm}|}{\textbf{Thèmes d’étude transversaux}} \\
		Modèles définis par une fonction d’une variable (optimisation, coût, etc.) & \casevide & \casevide \\
		Modèles d’évolution (discrets et continus, équations différentielles) & \casevide & \casevide \\
		Approche historique de la fonction logarithme & \casevide & \casevide \\
		Calculs d’aires (méthodes historiques, rectangles, Monte-Carlo) & \casevide & \casevide \\
		Répartition des richesses (courbe de Lorenz, indice de Gini) & \casevide & \casevide \\
		Inférence bayésienne (probabilités conditionnelles, formule de Bayes) & \casevide & \casevide \\
		Répétition d’expériences indépendantes, échantillonnage & \casevide & \casevide \\
		Temps d’attente (lois géométrique et exponentielle) & \casevide & \casevide \\
		Corrélation et causalité (régression, regard critique) & \casevide & \casevide \\
		\hline
		
		% --- Algorithmique et programmation ---
		\multicolumn{3}{|L{14cm}|}{\textbf{Algorithmique et programmation}} \\
		Écriture d’algorithmes (langage naturel ou Python) & \casevide & \casevide \\
		Boucles, tests, fonctions, listes & \casevide & \casevide \\
		Algorithmes spécifiques (dichotomie, Euler, Briggs, Monte-Carlo, etc.) & \casevide & \casevide \\
		\hline
	\end{longtable}
	
	\section{Calendrier et programmation harmonieuse}
	
	\subsection{Calendrier retenu (zone C, jours ouvrés)}
	Les vacances et jours fériés ont été pris en compte. Les semaines de cours sont les suivantes :
	
	\begin{center}
		\begin{tabular}{|C{3.5cm}|C{5cm}|C{2.5cm}|}
			\hline
			\textbf{Période} & \textbf{Dates} & \textbf{Nb de semaines} \\
			\hline
			Rentrée – Toussaint & 1er sept. – 16 oct. 2026 & 7 \\
			Toussaint – Noël & 2 nov. – 18 déc. 2026 & 7 \\
			Noël – Hiver & 4 janv. – 5 fév. 2027 & 5 \\
			Hiver – Printemps & 22 fév. – 2 avr. 2027 & 6 \\
			Printemps – Fin mai & 19 avr. – 31 mai 2027 & 6 \\
			\hline
			\textbf{Total} & & \textbf{31 semaines} \\
			\hline
		\end{tabular}
	\end{center}
	
	\textit{Jours fériés (tombant un jour de cours) exclus :} 1er janvier 2027 (vendredi), 29 mars 2027 (lundi de Pâques), 6 mai 2027 (Ascension, jeudi), 17 mai 2027 (Pentecôte, lundi). Les 1er et 8 mai 2027 tombent un samedi.
	
	\subsection{Programmation par trimestre (approche spiralée)}
	Les tableaux suivants précisent pour chaque semaine :
	\begin{itemize}
		\item les contenus nouveaux,
		\item les \textbf{points ciblés dans l’évaluation formative} (5 à 10 min en début de séance) incluant systématiquement une question Python,
		\item des remarques utiles.
	\end{itemize}
	Les évaluations formatives mêlent mathématiques et programmation, avec réactivation des notions antérieures.
	
	% --- 1er trimestre ---
	\newpage
	\subsubsection{1\textsuperscript{er} trimestre (septembre – novembre 2026) – 14 semaines}
	
	\rowcolors{2}{gray!8}{white}
	\begin{longtable}{|C{0.8cm}|L{4.0cm}|L{5.8cm}|L{2.5cm}|}
		\hline
		\textbf{Sem.} & \textbf{Contenus nouveaux} & \textbf{Évaluation formative -- points évalués (dont Python)} & \textbf{Remarques} \\
		\hline
		\endhead
		\hline
		\multicolumn{4}{r}{Suite de la page suivante} \\
		\endfoot
		\hline
		\endlastfoot
		
		1-2 & Suites : limites finies/infinies, opérations, théorème des gendarmes. & Calculs de puissances/racines. Suites arithmétiques/géométriques (1\textsuperscript{re}). \textbf{Python} : boucle \texttt{for} pour calculer les premiers termes d’une suite. & Introduction des modèles discrets. \\
		\hline
		3-4 & Suites géométriques (sommes, limite). Suites arithmético-géométriques. & \textbf{Spirale :} limite de suite géométrique (S1-2). Équations du 1er degré. \textbf{Python} : programme de recherche de seuil pour une suite. & Algorithmes : seuils, calcul de termes. \\
		\hline
		5-6 & Continuité, TVI, réciproque d’une fonction monotone. & \textbf{Spirale :} limites de suites (S3-4). Dérivées des fonctions de référence. \textbf{Python} : implémentation de la dichotomie. & Lien suites/fonctions (points fixes). \\
		\hline
		7-8 & Logarithme népérien (déf., équation fonctionnelle, dérivée). & \textbf{Spirale :} TVI (S5-6). Dérivée de l’exponentielle. \textbf{Python} : algorithme de Briggs pour tabuler des logarithmes. & Approche historique (Neper, Briggs). \\
		\hline
		9-10 & Dérivée de $f(ax+b)$, $e^{u(x)}$, $\ln u(x)$, $u(x)^2$. & \textbf{Spirale :} propriétés du log (S7-8). Calculs de limites. \textbf{Python} : représentation graphique de fonctions avec \texttt{matplotlib}. & Optimisation. \\
		\hline
		11-12 & Primitives, équation différentielle $y'=ay+b$. & \textbf{Spirale :} dérivées composées (S9-10). Limites de suites. \textbf{Python} : méthode d’Euler pour l’approximation de solutions. & Modèles d’évolution. \\
		\hline
		13-14 & Thème \emph{Modèles d’évolution} (discret / continu). & \textbf{Bilan partiel T1 :} suites, limites, log, dérivées, primitives, équadiff. \textbf{Python} : comparaison Euler / solution exacte. & Synthèse + Euler. \\
		\hline
		\multicolumn{4}{l}{} \\
		\hline
		\multicolumn{4}{l}{\textbf{Contrôles continus :} semaines 5, 10, 14 (avec exercice Python). \textbf{Formatives :} chaque semaine.} \\
		\hline
	\end{longtable}
	
	% --- 2e trimestre ---
	\newpage
	\subsubsection{2\textsuperscript{e} trimestre (novembre 2026 – février 2027) – 12 semaines}
	
	\rowcolors{2}{gray!8}{white}
	\begin{longtable}{|C{0.8cm}|L{4.0cm}|L{5.8cm}|L{2.5cm}|}
		\hline
		\textbf{Sem.} & \textbf{Contenus nouveaux} & \textbf{Évaluation formative -- points évalués (dont Python)} & \textbf{Remarques} \\
		\hline
		\endhead
		\hline
		\multicolumn{4}{r}{Suite de la page suivante} \\
		\endfoot
		\hline
		\endlastfoot
		
		15-16 & Fonctions convexes (dérivée seconde, caractérisation, point d’inflexion). & \textbf{Spirale :} étude complète d’une fonction (T1). Calcul de dérivées secondes. \textbf{Python} : tracé de courbes et de tangentes pour visualiser la convexité. & Applications aux courbes. \\
		\hline
		17-18 & Intégrale d’une fonction positive, relation de Chasles, valeur moyenne. & \textbf{Spirale :} convexité (S15-16). Primitives immédiates (S11-12). \textbf{Python} : méthode des rectangles (sommes de Riemann). & Méthode des rectangles. \\
		\hline
		19-20 & Théorème fondamental, calcul d’intégrales par primitives. & \textbf{Spirale :} rectangles (S17-18). Suites. \textbf{Python} : méthode de Monte-Carlo pour estimer une aire. & \textbf{PBL 1 (2h)} : Quadrature de la parabole (avec Python). \\
		\hline
		21-22 & Lois discrètes : Bernoulli, binomiale, géométrique. & \textbf{Spirale :} calcul intégral (S19-20). Logarithme (S7-8). \textbf{Python} : simulation de tirages de Bernoulli et calcul de fréquences. & Simulations Python. \\
		\hline
		23-24 & Thème \emph{Répétition d’expériences, échantillonnage}. & \textbf{Spirale :} loi binomiale (espérance, écart-type) (S21-22). Fonctions (limites). \textbf{Python} : génération d’échantillons, calcul d’intervalles de fluctuation. & Intervalles de fluctuation. \\
		\hline
		25-26 & Statistique à deux variables (nuage, régression, corrélation). & \textbf{Spirale :} échantillonnage (S23-24). Convexité (S15-16). \textbf{Python} : calcul du point moyen, tracé du nuage et de la droite des moindres carrés. & \textbf{PBL 2 (2h)} : Corrélation et causalité (avec Python). \\
		\hline
		\multicolumn{4}{l}{} \\
		\hline
		\multicolumn{4}{l}{\textbf{Contrôles continus :} semaines 18, 22, 26 (avec exercice Python). \textbf{Formatives :} chaque semaine.} \\
		\hline
	\end{longtable}
	
	% --- 3e trimestre ---
	\newpage
	\subsubsection{3\textsuperscript{e} trimestre (février – mai 2027) – 12 semaines}
	
	\rowcolors{2}{gray!8}{white}
	\begin{longtable}{|C{0.8cm}|L{4.0cm}|L{5.8cm}|L{2.5cm}|}
		\hline
		\textbf{Sem.} & \textbf{Contenus nouveaux} & \textbf{Évaluation formative -- points évalués (dont Python)} & \textbf{Remarques} \\
		\hline
		\endhead
		\hline
		\multicolumn{4}{r}{Suite de la page suivante} \\
		\endfoot
		\hline
		\endlastfoot
		
		27-28 & Lois à densité (uniforme, exponentielle). & \textbf{Spirale :} loi géométrique (S21-22). Intégration (aires). \textbf{Python} : simulation d’une loi exponentielle à partir de la loi uniforme. & Absence de mémoire. \\
		\hline
		29-30 & Thème \emph{Temps d’attente} (géométrique vs exponentielle). & \textbf{Spirale :} loi exponentielle (S27-28). Suites arithmético-géométriques (S3-4). \textbf{Python} : calcul de demi-vie par simulation sur un grand échantillon. & \textbf{PBL 3 (2h)} : Demi-vie d’un radioélément (avec Python). \\
		\hline
		31-32 & Inférence bayésienne (probas conditionnelles, formule de Bayes). & \textbf{Spirale :} probas conditionnelles (1\textsuperscript{re}). Loi exponentielle (S27-28). \textbf{Python} : simulation de tests de dépistage (vrais/faux positifs). & Tests de dépistage. \\
		\hline
		33-34 & Thème \emph{Répartition des richesses} (Lorenz, Gini). & \textbf{Spirale :} convexité (S15-16). Calcul intégral (S19-20). \textbf{Python} : calcul numérique de l’indice de Gini à partir de données. & Applications socio-économiques. \\
		\hline
		35-36 & Synthèse : modèles par fonction d’une variable. & \textbf{Bilan analyse :} dérivées, intégrales, suites. \textbf{Python} : optimisation d’une fonction par recherche de seuil / dichotomie. & Optimisation, coûts. \\
		\hline
		37-38 & Révisions et évaluations finales. & \textbf{Bilan probabilités + Python :} toutes les simulations et algorithmes vus dans l’année. & Préparation aux épreuves. \\
		\hline
		\multicolumn{4}{l}{} \\
		\hline
		\multicolumn{4}{l}{\textbf{Contrôles continus :} semaines 30, 34, 38 (avec exercice Python). \textbf{Formatives :} chaque semaine.} \\
		\hline
	\end{longtable}
	
	\section{Évaluations}
	
	\subsection{Évaluations formatives (hebdomadaires)}
	Chaque semaine, une courte activité (5–10 min) mêle mathématiques et Python. Les points évalués sont détaillés dans les tableaux ci-dessus. Ces évaluations ne sont pas notées, mais elles permettent une auto-évaluation et un suivi des compétences algorithmiques.
	
	\subsection{Évaluations sommatives (contrôles continus)}
	Au moins trois contrôles par trimestre, d’une durée de 1 h chacun.  
	Chaque contrôle comporte :
	\begin{itemize}
		\item une partie sans calculatrice (calculs, raisonnement),
		\item une partie avec calculatrice / Python (exercice de programmation : écriture de fonction, complétion de code, simulation, représentation graphique, ou recherche de seuil).
	\end{itemize}
	
	\begin{center}
		\begin{tabular}{|C{3.5cm}|C{3.5cm}|C{3.5cm}|C{3.5cm}|}
			\hline
			\textbf{Trimestre} & \textbf{Contrôle 1} & \textbf{Contrôle 2} & \textbf{Contrôle 3} \\
			\hline
			1\textsuperscript{er} (sem. 5, 10, 14) & Suites, limites, TVI \newline + Python (boucles, seuils) & Logarithme, dérivées \newline + Python (Briggs, graphiques) & Primitives, équa. diff. \newline + Python (Euler) \\
			\hline
			2\textsuperscript{e} (sem. 18, 22, 26) & Convexité, intégration (aires) \newline + Python (rectangles, Monte-Carlo) & Lois binomiale et géométrique \newline + Python (simulations) & Statistique à 2 variables \newline + Python (régression) \\
			\hline
			3\textsuperscript{e} (sem. 30, 34, 38) & Lois à densité, exponentielle \newline + Python (simulation, demi-vie) & Inférence bayésienne \newline + Python (tests de dépistage) & Synthèse (modèles, optimisation) \newline + Python (projet complet) \\
			\hline
		\end{tabular}
	\end{center}
	
	\subsection{Problèmes basés sur l’apprentissage (PBL) – 2 h chacun}
	Trois PBL sont répartis sur l’année, en séances de 2 h (travail en groupes, compte rendu écrit et code Python rendu) :
	
	\begin{enumerate}
		\item \textbf{PBL 1 (semaine 20)} : \emph{Quadrature de la parabole}. Approcher une aire par la méthode d’Archimède, puis par les rectangles, et comparer avec le calcul intégral. \textbf{Python} : implémentation des méthodes et comparaison graphique.
		\item \textbf{PBL 2 (semaine 26)} : \emph{Corrélation et causalité}. À partir de données réelles (température/CO$_2$), ajuster une droite de régression, discuter de la corrélation et des pièges de l’extrapolation. \textbf{Python} : calculs statistiques et tracé du nuage.
		\item \textbf{PBL 3 (semaine 30)} : \emph{Demi-vie d’un radioélément}. Modéliser la désintégration par un modèle discret (suite) et continu (exponentielle), simuler avec Python, déterminer la demi-vie.
	\end{enumerate}
	
	\section{Devoirs maison (DM) -- progression, Python et utilisation de l’IA}
	
	Afin de favoriser une acquisition à long terme, un devoir maison est distribué \textbf{toutes les deux semaines} (18 DM sur l’année).  
	Chaque DM porte sur l’ensemble du programme traité depuis le 1er septembre, et comporte :
	\begin{itemize}
		\item des exercices de calcul et de raisonnement,
		\item un exercice de programmation Python,
		\item une mission spécifique d’utilisation de l’IA (consultation, débuggage, confrontation de méthodes ou critique de réponse).
	\end{itemize}
	
	L’élève doit joindre à sa copie un \textbf{journal de bord} retraçant ses échanges avec l’IA et les décisions prises à la suite de ces échanges.
	
	\begin{center}
		\small
		\rowcolors{2}{gray!8}{white}
		\begin{longtable}{|C{0.8cm}|C{1cm}|L{6.8cm}|L{6.8cm}|}
			\hline
			\textbf{DM} & \textbf{Sem.} & \textbf{Thèmes évalués (mathématiques + Python)} & \textbf{Mission IA spécifique} \\
			\hline
			\endhead
			\hline
			\multicolumn{4}{r}{Suite à la page suivante} \\
			\endfoot
			\hline
			\endlastfoot
			
			1 & 3 & Suites (limites, géométriques). Python : calcul des termes et graphique. & Demander à l’IA d’expliquer la notion de convergence avec un exemple ; comparer avec la définition du cours. \\
			\hline
			2 & 5 & Limites de suites, TVI. Python : dichotomie. & Soumettre à l’IA un script bogué de dichotomie, lui demander de l’analyser, puis justifier les corrections. \\
			\hline
			3 & 7 & Suites arithmético-géométriques, logarithme (équation fonctionnelle). Python : algorithme de Briggs. & Demander à l’IA de proposer une méthode historique alternative pour le log, puis discuter des limites pratiques. \\
			\hline
			4 & 9 & Logarithme (dérivée, limites), fonctions composées. Python : tracé de courbes et recherche de seuil. & Utiliser l’IA pour générer 3 fonctions composées différentes, les étudier et vérifier les résultats. \\
			\hline
			5 & 11 & Primitives, équations différentielles $y'=ay+b$. Python : méthode d’Euler. & Demander à l’IA de résoudre l’équation exacte, puis comparer avec l’approximation d’Euler sur un intervalle donné. \\
			\hline
			6 & 13 & Modèles d’évolution (discret/continu). Python : comparaison Euler / exacte. & Faire générer par l’IA un scénario concret (ex. température) et modéliser les deux approches. \\
			7 & 16 & Convexité (dérivée seconde, inflexion). Python : tracé des tangentes. & Demander à l’IA d’expliquer la convexité via la position des cordes, puis illustrer avec un code Python. \\
			\hline
			8 & 18 & Intégration (aire, valeur moyenne). Python : méthode des rectangles. & Utiliser l’IA pour générer une fonction aléatoire, estimer son intégrale par rectangles, puis vérifier. \\
			\hline
			9 & 20 & Calcul intégral (théorème fondamental). Python : Monte-Carlo. & Confronter l’IA sur la convergence de Monte-Carlo (loi des grands nombres) et interpréter sa réponse. \\
			\hline
			10 & 22 & Lois discrètes (Bernoulli, binomiale, géométrique). Python : simulations. & Demander à l’IA de simuler 1000 lancers de pièce en Python, comparer avec les résultats théoriques. \\
			\hline
			11 & 24 & Échantillonnage, intervalles de fluctuation. Python : génération d’échantillons. & Utiliser l’IA pour discuter de la représentativité d’un échantillon et proposer des tests de validation. \\
			\hline
			12 & 26 & Statistique à deux variables (régression). Python : point moyen, droite des MC. & Faire générer par l’IA un nuage de points et comparer plusieurs méthodes d’ajustement. \\
			\hline
			13 & 28 & Lois à densité (uniforme, exponentielle). Python : simulation exponentielle. & Demander à l’IA de démontrer la propriété d’absence de mémoire, puis la vérifier par simulation. \\
			\hline
			14 & 30 & Temps d’attente (géométrique vs exponentielle). Python : calcul de demi-vie. & Utiliser l’IA pour déterminer la demi-vie du Carbone 14 et confronter le résultat avec les données. \\
			\hline
			15 & 32 & Inférence bayésienne (probas conditionnelles). Python : tests de dépistage. & Demander à l’IA de calculer les valeurs prédictives pour une maladie rare, interpréter et critiquer. \\
			\hline
			16 & 34 & Répartition des richesses (Lorenz, Gini). Python : calcul de l’indice. & Faire générer par l’IA des données de répartition, calculer l’indice de Gini et commenter. \\
			\hline
			17 & 36 & Synthèse : optimisation par fonction d’une variable. Python : balayage / dichotomie. & Utiliser l’IA pour proposer un problème d’optimisation économique original, puis le résoudre. \\
			\hline
			18 & 38 & Bilan général (tout le programme). Python : projet complet. & Demander à l’IA un retour sur la structure du code (lisibilité, efficacité) et améliorer le script. \\
			\hline
		\end{longtable}
	\end{center}
	
	\textbf{Remarques :}
	\begin{itemize}
		\item Les DM sont à rendre une semaine après la distribution.
		\item La mission IA doit être clairement identifiée dans le journal de bord (copie d’écran ou retranscription des questions/réponses, suivie de l’analyse personnelle de l’élève).
		\item La correction est effectuée en classe, avec une attention particulière aux erreurs fréquentes, aux méthodes de résolution et aux bonnes pratiques d’utilisation de l’IA.
	\end{itemize}
	
	\section{Conseils pour la mise en œuvre}
	\begin{itemize}
		\item \textbf{Rituels} : consacrer 5 à 10 minutes par séance à des exercices de calcul (mental, littéral) et à de courts extraits de code Python (lecture, complétion, débogage).
		\item \textbf{Travail en îlots} : favoriser les échanges entre élèves lors des phases de recherche et des PBL.
		\item \textbf{Numérique} : utiliser Python régulièrement (au moins une séance sur deux). Insister sur la programmation modulaire (fonctions) et la documentation du code.
		\item \textbf{IA éthique} : former les élèves à formuler des requêtes précises et à évaluer de manière critique les réponses de l’IA (vérification des calculs, cohérence, limites). Insister sur le fait que l’IA est un \emph{outil d’aide}, pas un substitut au raisonnement personnel.
		\item \textbf{Différenciation} : proposer des exercices de difficulté progressive et des approfondissements pour les élèves plus à l’aise (notamment en programmation et en interaction avec l’IA).
		\item \textbf{Trace écrite} : veiller à ce que chaque élève dispose d’un cahier de cours structuré (définitions, propriétés, exemples) et d’un cahier de programmation regroupant les fonctions et scripts clés.
	\end{itemize}
	
	\section*{Conclusion}
	Cette programmation respecte les 31 semaines de cours disponibles et couvre l’intégralité du programme, avec une intégration systématique de Python et de l’utilisation critique de l’IA dans les apprentissages et les évaluations. La progressivité spiralaire, appliquée aussi aux compétences algorithmiques et numériques, garantit que les notions mathématiques, la programmation et le jugement critique sont solidement maîtrisés en fin d’année. Les évaluations variées (formatives, sommatives, PBL, DM) permettent de suivre la progression de chaque élève et de développer ses compétences de modélisation, de raisonnement, de communication, de codage et de collaboration avec les outils d’intelligence artificielle.
	
\end{document}